\section{Categorical QDM markers and one stabilization}
\label{sec:atomic-one-step}

We prove the unconditional theorem by applying the quantum \(D\)-module
decomposition formulas directly.  The common argument is stated first for an
arbitrary additive marker.  The rest of the section constructs the smallest
marker needed for a cubic threefold and verifies its center nullity.  The
framed invariant of Section~\ref{sec:framed-monodromy} will apply the same
ledger theorem to a block type carrying additional framing data.

\subsection{The one-blowup model}

Let \(Y\) be smooth and projective.  Pass to numerical Novikov variables,
restrict the big quantum connection to even bulk variables, and then pass to
its generic coefficient field.  After an algebraic extension, multiplication
by the Euler field has a primary decomposition.  We call the resulting
connection summands the \emph{generic even QDM blocks} of \(Y\).  A block
retains its connection, grading, pairing, and Euler endomorphism.  A block type
may require additional marked data.  Its morphisms preserve all the typed data
and are regular connection isomorphisms after scalar extension and formal
change of bulk coordinates.  In particular, their gauges use neither negative
nor fractional powers of the connection variable \(z\).

Fix a block type \(\mathcal T\).  Let \(\Pi_{\mathcal T}\) be the
regular-isomorphism groupoid of blocks of that type and put
\[
 \mathcal L_{\mathcal T}
 :=\pi_0\operatorname{Sym}^{\sqcup}(\Pi_{\mathcal T}),
\]
the effective free commutative monoid on their isomorphism classes.  Direct
sum gives each \(Y\) a block ledger
\(\mathcal B_{\mathcal T}(Y)\in\mathcal L_{\mathcal T}\).  A
\emph{marker fold} is a commutative-monoid homomorphism
\[
 w:\mathcal L_{\mathcal T}\longrightarrow A,
 \qquad I_w(Y):=w(\mathcal B_{\mathcal T}(Y)),
\]
where \(A\) is an effective commutative monoid.  No group completion is
made: it can destroy useful folds, including Boolean ones.

The two specializations use the following notation.
\begin{table}[ht]
\centering
\begin{tabular}{@{}ll@{}}
\toprule
Symbol & Meaning \\
\midrule
\(\Pi_{\mathcal T}\) & regular-isomorphism groupoid of typed QDM blocks \\
\(\mathcal L_{\mathcal T}\) &
  \(\pi_0\operatorname{Sym}^{\sqcup}(\Pi_{\mathcal T})\), the effective block monoid \\
\(\mathcal B_{\mathcal T}(Y)\) & block ledger of \(Y\) \\
\(w:\mathcal L_{\mathcal T}\to A\) & additive marker fold into an effective monoid \\
\(\omega_j\) & the \(j\)-th center occurrence in one blowup \\
\(c_w(\omega_j)\) & marker of that actual center occurrence \\
\bottomrule
\end{tabular}
\caption{Notation for the occurrence-indexed marker ledger.}
\label{tab:marker-notation}
\end{table}

Consider one blowup
\(\beta:\widetilde Y=\Bl_ZY\to Y\) along a smooth center of codimension
\(c\ge2\).  Iritani's comparison has one ambient summand and \(c-1\) center
summands.  We retain the occurrence data of each center summand: for the
\(j\)-th, record
\[
 \omega_j=(\beta,Z,j,\varsigma_j,\chi_j),
 \qquad 0\le j\le c-2,
\]
where \(\varsigma_j\) is Iritani's center reconstruction parameter and
\(\chi_j\) is the induced numerical Novikov specialization.  Its correction
is
\(c_w(\omega_j):=w(\mathcal B_{\mathcal T}(Z;\omega_j))\).  Equality is
required only after applying \(w\); distinct block multisets may have the same
marker.

\begin{theorem}[Occurrence-indexed marker ledger]
\label{thm:marker-ledger}
\coverage{fragment}%
\lean{packet_multiplicity_eq_of_preserving_chain,%
  packet_multiplicity_eq_of_typed_weak_factorization}%
\imports{AKMW:weak-factorization}%
Let \(w:\mathcal L_{\mathcal T}\to A\) be a marker fold unchanged by the
allowed regular scalar extensions and formal coordinate changes.  Suppose
\begin{align}
 I_w(\PP_Y(V))&=rI_w(Y),
 &&\rk V=r,                                      \label{eq:marker-pbundle}\\
 I_w(\Bl_ZY)&=I_w(Y)+\sum_{j=0}^{c-2}c_w(\omega_j),
 &&\operatorname{codim}_YZ=c\ge2.                \label{eq:marker-blowup}
\end{align}
If every actual smooth center occurrence of dimension at most \(d-2\) has
zero correction, then \(I_w\) is a birational invariant of smooth projective
\(d\)-folds.
\end{theorem}

\begin{proof}
Projective weak factorization writes a birational map of smooth projective
\(d\)-folds as a zigzag of blowups and blowdowns along smooth centers of
codimension at least two \cite[Theorem~0.1.1]{AKMW}.  By
\eqref{eq:marker-blowup} and center nullity, the marker is unchanged at each
step; the equalities therefore telescope.  Disconnected centers are treated
componentwise, while a divisor blowup is an isomorphism.
\end{proof}

Equivalently, the marker descends on object sets:
\begin{equation}
\begin{CD}
 \operatorname{Ob}(\mathsf{SmProj}_d) @>{I_w}>> A\\
 @VVV @|\\
 \operatorname{Ob}(\mathsf{SmProj}_d)/{\sim_{\mathrm{bir}}}
 @>>{\overline I_w}> A.
\end{CD}
\label{diag:marker-descent}
\end{equation}
This square reformulates the proof; it is not an additional input.

\subsection{The QDM operation providers}

Three adapters place the comparison theorems in the preceding interface.
First, the intrinsic ample-adic completion and the Laurent neighborhood of a
fibre or exceptional cusp need not map to one another.  Both receive an
injective map from the algebraic coefficient algebra on which the finite
connection matrix is defined.  Its fraction field is a common generic spine.
Rank, primary factors, nilpotence, regular formal gauges, and conjugacy
invariants are unchanged after either field extension.

Second, the comparisons are regular in \(z\).  Iritani's Theorem~5.18 and
Iritani--Koto's Theorem~5.1 are respectively defined over rings of the form
\[
 \C[z]((q_{\mathrm{exc}}^{-1/s}))[[Q,\widetilde\tau]],
 \qquad
 \C[z]((q_{\mathrm{fib}}^{-1/r'}))[[Q,\widehat\tau]],
\]
where Iritani--Koto take \(r'=r\) when \(r-1\) is even and \(r'=2r\) when
\(r-1\) is odd.  Their inverses use no negative powers of \(z\); see also
\cite[Remark~5.3]{IritaniKoto}.  Homogeneity and the explicit pull-push and
Fourier formulas preserve parity after odd variables are set to zero.  The
combined bulk Jacobians are invertible, hence so are their even-even blocks.

Third, Iritani's center map need not be injective
\cite[(5.15)]{IritaniBlowup}.  After numerical reduction it has the form
\begin{equation}
 Q_Z^d\longmapsto
 Q^{i_*d}q_{\mathrm{exc}}^{-\rho_Z\cdot d/(c-1)}.
 \label{eq:center-novikov-map-atomic}
\end{equation}
Choose divisors \(D_1,\ldots,D_\rho\) separating \(N_1(Z)\) and retain the
characters \(\exp(\sum_k(D_k\cdot d)s_k)\).  Every fibre of
\eqref{eq:center-novikov-map-atomic} is finite: an ample degree from \(Y\)
bounds the lattice points in the corresponding slice of
\(\overline{NE}(Z)\).  On a fibre of size \(n\), a generic one-parameter
restriction gives distinct exponentials.  Their derivatives of orders
\(0,\ldots,n-1\) form an invertible Vandermonde matrix, so the
character-tagged map is injective.
Remarks~2.3 and 5.6 of \cite{IritaniBlowup} show that divisor and Novikov
variables occur in precisely these combinations.  Each center summand is a
faithful scalar extension and formal reparametrization of the intrinsic
generic numerical QDM of \(Z\).

\begin{proposition}[QDM operation ledgers]
\label{prop:qdm-operation-ledgers}
\coverage{absent}%
\imports{IritaniKoto:thm-5.1,IritaniBlowup:thm-5.18}%
Let a marker fold be preserved by the regular scalar extensions and formal
coordinate changes described above.  The projective-bundle and blowup QDM
comparisons then supply \eqref{eq:marker-pbundle} and
\eqref{eq:marker-blowup}; each center marker is evaluated at its actual
occurrence \(\omega_j\).
\end{proposition}

\begin{proof}
For a rank-\(r\) bundle, Iritani--Koto's Theorem~5.1 gives \(r\) copies of
the base QDM after Laurent extension \cite[Theorem~5.1]{IritaniKoto}.  Its
invertible bulk Jacobian separates their unit coordinates.  The regularity,
parity, and generic-spine adapters carry the marker across the comparison,
and additivity gives \eqref{eq:marker-pbundle}.  Tensoring the bundle by a
line bundle, when needed for the global-generation convention, changes
neither the projective bundle nor the intrinsic Novikov embedding
\cite[Remark~1.2 and Remark~5.2]{IritaniKoto}.

For a codimension-\(c\) blowup, Iritani's Theorem~5.18 gives one ambient QDM
and \(c-1\) center QDMs \cite[Theorem~5.18]{IritaniBlowup}.  Parts (4), (6),
and (7) give separated leading branches and an invertible combined bulk
Jacobian.  Apply the same regularity and parity adapters, and apply the
faithful character adapter to each center copy.  Additivity then gives
\eqref{eq:marker-blowup}.  Each comparison is evaluated over its own
coefficient spine; the proof never composes Laurent comparisons recursively.
\end{proof}

The QDM comparison need not identify the two ledger elements before applying
the fold.  The relevant categorical operation is instead the additivity square
\begin{equation}
\begin{CD}
 \mathcal L_{\mathcal T}\times\mathcal L_{\mathcal T}
 @>{+}>> \mathcal L_{\mathcal T}\\
 @V{w\times w}VV @VV{w}V\\
 A\times A @>>{+}> A.
\end{CD}
\label{diag:marker-fold-additivity}
\end{equation}
Proposition~\ref{prop:qdm-operation-ledgers} supplies only the resulting
equalities in \(A\), as required by Theorem~\ref{thm:marker-ledger}.

\subsection{The rank-two residue marker}

Let \(\mathcal T_{\mathrm{at}}\) be the block type retaining the connection,
grading, pairing, and Euler endomorphism, with no marked bulk section or loop
framing, and specialize to \(A=\mathbf N\).  Let \(E\) be a generic even
block with centered connection
\begin{equation}
 z\partial_z y=A(z)y,
 \qquad A(z)=z^{-1}N+A_0+zA_1+z^2A_2+\cdots.
 \label{eq:atomic-block-expansion}
\end{equation}
Assume \(\rk E=2\), \(N\ne0\), and \(N^2=0\).  Then
\(L=\im N=\ker N\) is a line.  Define
\begin{equation}
 E^\sharp=\{s\in E:s\bmod z\in L\}.
 \label{eq:atomic-modification}
\end{equation}

\begin{lemma}[The regular coefficient preserves the nilpotent line]
\label{lem:A0preserve}
\coverage{conditional_deduction}%
\lean{atomicRankTwo_regularCoefficient_preserves_nilpotentLine,%
  atomicRankTwo_pairingEquations_of_loopHorizontality}%
The Poincar\'e pairing implies \(A_0L\subset L\).
\end{lemma}

\begin{proof}
Separated spectral factors are orthogonal for the horizontal Poincar\'e
pairing: coefficientwise horizontality kills every off-diagonal block by an
invertible Sylvester operator.  Hence the restriction to \(E\) is
nondegenerate.  Write it as \(P(z)=P_0+zP_1+\cdots\).  Horizontality gives
\[
 z\partial_zP(z)+A(z)^TP(z)+P(z)A(-z)=0.
\]
The \(z^{-1}\) coefficient gives \(N^TP_0=P_0N\).  Thus \(L\) is isotropic
and, in dimension two, \(L^\perp=L\).  The constant coefficient is
\[
 A_0^TP_0+P_0A_0+N^TP_1-P_1N=0.
\]
Evaluating this identity on \((x,x)\), where \(x\in L=\ker N\), gives
\((A_0x,x)=0\).  Hence \(A_0x\in L^\perp=L\).
\end{proof}

Choose a frame with \(L=\langle e_1\rangle\) and
\(Ne_2=\nu e_1\), \(\nu\ne0\).  The modified lattice has frame
\(e_1,ze_2\), so
\begin{equation}
 A^\sharp=S^{-1}AS-S^{-1}z\partial_zS,
 \qquad S=\diag(1,z).
 \label{eq:Asharp}
\end{equation}
The term \(z^{-1}N\) becomes regular.  The only possible new pole is the
\(E_{21}\)-coefficient of \(A_0\), which vanishes by
Lemma~\ref{lem:A0preserve}.  Thus
\(A^\sharp=R+zR_1+\cdots\).  Define
\begin{equation}
 \delta^\sharp(E):=(\tr R)^2-4\det R.
 \label{eq:res-disc}
\end{equation}

\begin{proposition}[Rank-two rigidity]
\label{prop:rank2-rigidity}
\coverage{conditional_deduction}%
\lean{atomicRankTwo_modifiedBase_regular_of_flat_connection,%
  atomicRankTwo_modifiedResidue_lax_of_flat_connection,%
  atomicRankTwo_residueDiscriminant_frozen_by_horizontality_and_flatness,%
  atomicRankTwo_residueDiscriminant_constant_over_formal_germ,%
  atomicRankTwo_modifiedBase_pole_eq_zero,%
  atomicRankTwo_residueDiscriminant_constant_along_base}%
Under formal change of even bulk coordinates, the modified connection is
regular in every base direction and
\begin{equation}
 \partial R=[G_\partial,R]
 \label{eq:atomic-lax}
\end{equation}
for a regular matrix \(G_\partial\).  Consequently
\(\partial(\delta^\sharp)=0\).  Regular connection isomorphisms preserve the
domain and value of \(\delta^\sharp\).
\end{proposition}

\begin{proof}
Write a centered base equation as
\(\partial y=(z^{-1}C_\partial+C_{\partial,0}+O(z))y\).  The
\(z^{-2}\) coefficient of flatness gives \([N,C_\partial]=0\), and the trace
of the \(z^{-1}\) coefficient gives \(\tr C_\partial=0\); here the base
equation has been scalar-centered.  The commutant of a nonzero rank-one
nilpotent \(2\times2\) matrix is the coefficient-field span of \(I\) and
\(N\), so \(C_\partial=q_\partial N\).  The same \(z^{-1}\) coefficient gives
\[
 \partial N=-q_\partial N+
 [N,q_\partial A_0-C_{\partial,0}].
\]
Thus formal bulk transport preserves the nonzero square-zero orbit and its
image line.  After modification, the only possible base
pole is \(z^{-1}kE_{21}\).  If
\(R=\left(\begin{smallmatrix}a&\nu\\c&d\end{smallmatrix}\right)\), the
\(z^{-1}\) flatness equation is
\(kE_{21}+[R,kE_{21}]=0\).  Its diagonal entries are \(\nu k\) and
\(-\nu k\), so \(k=0\).  The constant term of the flatness equation is now
\eqref{eq:atomic-lax}; trace and determinant are constant under infinitesimal
conjugation.  A regular isomorphism carries \(N\), \(L\), and the modified
lattice to their counterparts and conjugates the residue.
\end{proof}

\begin{proposition}[Exponent interpretation]
\label{prop:residue-discriminant-exponents}
\coverage{fragment}%
\lean{residueDiscriminant_eq_squared_eigenvalue_separation}%
\imports{PutSinger:ch-3}%
If \(\rho_1,\rho_2\) are the eigenvalues of \(R\), their classes modulo
\(\Z\) are the formal exponent classes of the original centered block and
\[
 \delta^\sharp=(\rho_1-\rho_2)^2.
\]
\end{proposition}

\begin{proof}
The elementary modification uses integral powers of \(z\), so it preserves
formal exponent classes modulo \(\Z\).  The modified connection is regular
singular with residue \(R\); the regular-singular classification identifies
the exponent classes with the residue eigenvalues modulo \(\Z\)
\cite[Chapter~3]{PutSinger}.  The displayed identity is the quadratic
discriminant formula.
\end{proof}

Define the fold on each block by
\begin{equation}
 w_{\mathrm{at}}(E)=
 \begin{cases}
 1,&\rk E=2,\ N\ne0,\ N^2=0,\ \delta^\sharp(E)\ne0,\\
 0,&\text{otherwise},
 \end{cases}
 \label{eq:atomic-fold}
\end{equation}
and extend it additively to
\(w_{\mathrm{at}}:\mathcal L_{\mathcal T_{\mathrm{at}}}\to\mathbf N\).  Put
\(I_{\mathrm{at}}(Y)=
w_{\mathrm{at}}(\mathcal B_{\mathcal T_{\mathrm{at}}}(Y))\).
Proposition~\ref{prop:rank2-rigidity} makes this a lawful marker for
Proposition~\ref{prop:qdm-operation-ledgers}.

\subsection{The cubic block}
\label{subsec:shared-cubic-packet}

Let \(X\subset\PP^4\) be a smooth cubic threefold, let \(P=H|_X\), and put
\(q=Q^\ell\) for the line class and \(r=(3q)^{1/2}\).  Beauville's formulas
give \cite[main theorem and (2.1)--(2.3)]{Beauville}
\begin{equation}
 P\star P=P^2+6q,\qquad
 P\star P^2=P^3+15qP,\qquad
 P\star P^3=6qP^2+36q^2.
 \label{eq:beauville-cubic-products}
\end{equation}
Thus, in the basis \((1,P,P^2,P^3)\),
\begin{equation}
 K_X=2\begin{pmatrix}
 0&6q&0&36q^2\\1&0&15q&0\\0&1&0&6q\\0&0&1&0
 \end{pmatrix},
 \qquad G_X=\frac12\diag(3,1,-1,-3),
 \label{eq:cubic-small-even-connection}
\end{equation}
and \(z^2\partial_zS=(K_X+zG_X)S\).  This agrees with
\cite[Section~3]{Cai}.

\begin{proposition}[Small-even cubic block reduction]
\label{prop:cubic-block-data}
\coverage{fragment}%
\lean{cubicSmallEven_blockReduction,%
  cubicZeroBlock_modifiedResidue_indicialPolynomial,%
  cubicSmallEven_normalizedGauge_exists_and_unique,%
  cubicSmallEven_normalizedGauge_coefficients}%
\imports{Beauville:quantum-products}%
The characteristic and minimal polynomials of \(K_X\) are both
\(T^2(T^2-108q)\).  After adjoining \(r\), its eigenvalues are
\(6r,-6r,0,0\), and the zero generalized eigenspace is one rank-two Jordan
block.  A normalized gauge regular in \(z\) gives
\begin{equation}
 \begin{aligned}
 z^2\partial_z\widetilde S_+&=(6r+O(z^2))\widetilde S_+,\\
 z^2\partial_z\widetilde S_-&=(-6r+O(z^2))\widetilde S_-,\\
 z^2\partial_z
 \begin{pmatrix}\widetilde S_3\\\widetilde S_4\end{pmatrix}
 &=\bigl[J_0+zD_0+z^2E_0+O(z^3)\bigr]
 \begin{pmatrix}\widetilde S_3\\\widetilde S_4\end{pmatrix},
 \end{aligned}
 \label{eq:cubic-integral-block-reduction}
\end{equation}
where
\begin{equation}
 J_0=\begin{pmatrix}0&2\\0&0\end{pmatrix},\quad
 D_0=\begin{pmatrix}-19/18&0\\0&19/18\end{pmatrix},\quad
 E_0=\begin{pmatrix}0&-14/(81r^2)\\-8/81&0\end{pmatrix}.
 \label{eq:cubic-shared-block-data}
\end{equation}
\end{proposition}

\begin{proof}
The matrix
\[
 C=\begin{pmatrix}
 6r^3&-6r^3&0&-7r^2\\7r^2&7r^2&-2r^2&0\\
 3r&-3r&0&1\\1&1&1&0
 \end{pmatrix},
 \qquad \det C=-486r^5,
\]
puts \(K_X\) into
\[
 C^{-1}K_XC=
 \begin{pmatrix}6r&0&0&0\\0&-6r&0&0\\0&0&0&2\\0&0&0&0\end{pmatrix}.
\]
Partition as \(1\mid1\mid2\).  At each positive order in \(z\), the
off-diagonal gauge coefficient is the unique solution of a Sylvester
equation between blocks with disjoint scalar spectra.  The gauge and inverse
are regular in \(z\); explicitly, it has the normalized form
\begin{equation}
 A(z)=I+\sum_{n\ge1}A_nz^n,
 \qquad A_n\text{ block-off-diagonal}.
 \label{eq:cubic-integral-block-gauge}
\end{equation}
Direct multiplication through order \(z^2\) gives
\eqref{eq:cubic-shared-block-data}.
\end{proof}

For the zero block, \(N=J_0=2E_{12}\).  With \(S=\diag(1,z)\),
\begin{equation}
 R_X
 =2E_{12}+D_0-\frac8{81}E_{21}-\diag(0,1)
 =\begin{pmatrix}-19/18&2\\-8/81&1/18\end{pmatrix}.
 \label{eq:RX}
\end{equation}
Its characteristic polynomial is
\((T+1/6)(T+5/6)\), so
\begin{equation}
 \boxed{\delta^\sharp(X)=\frac49},
 \qquad I_{\mathrm{at}}(X)=1.
 \label{eq:cubic-delta}
\end{equation}

\subsection{Center nullity through dimension two}

\begin{proposition}[Low-dimensional nullity]
\label{prop:atomic-lowdim}
\coverage{conditional_deduction}%
\lean{cubicAtom_not_represented_by_curve,%
  cubicAtom_not_represented_by_surface}%
\uses{prop:rank2-rigidity}%
\imports{BeauvilleSurfaces:minimal-surface-classification}%
For every smooth projective \(Z\) of dimension at most two and every actual
QDM comparison occurrence \(\omega\),
\(I_{\mathrm{at}}(Z;\omega)=0\).
\end{proposition}

\begin{proof}
A point has one rank-one block.  For \(\PP^1\), the relation
\(p\star p=Qe^{t_p}\) gives two rank-one generic blocks.

Let \(C\) have genus \(g\ge1\).  There is no nonconstant genus-zero map to
\(C\), so its even big quantum product is classical.  With
\(\int_Cp=1\) and \(a=2-2g\), its centered even connection is
\[
 z\partial_zY=
 \left[z^{-1}aE_{21}+\diag(1/2,-1/2)\right]Y.
\]
For \(g=1\), \(N=0\).  For \(g>1\), the modification
\(S=\diag(z,1)\) gives residue \(aE_{21}-\tfrac12I\), of discriminant zero.
Thus every curve has zero marker.

Let \(S\) be a surface with \(K_S\) nef.  For a quantum-product term with
input degree \(b\), non-scalar Euler degree \(a\ge2\), curve class \(d\), and
nonunit even bulk degrees \(e_\nu\ge2\), the dimension axiom gives output
degree
\[
 a+b-2c_1(S)\cdot d+\sum_\nu(e_\nu-2)\ge b+2.
\]
The centered Euler operator strictly raises ordinary degree.  Its only block
is the whole even cohomology, of rank \(b_0+b_2+b_4\ge3\), so its marker is
zero.

A minimal surface with non-nef canonical class is \(\PP^2\) or a ruled
surface \(\PP_C(V)\) \cite[Chapter~VI]{BeauvilleSurfaces}.  The relation
\(H^{\star3}=Q\) makes \(\PP^2\) generically semisimple.  Formula
\eqref{eq:marker-pbundle} and the curve case give zero for ruled surfaces.
Every nonminimal surface is obtained from a minimal one by point blowups;
Formula~\eqref{eq:marker-blowup} adds only zero point terms.  This proves
intrinsic nullity.  The faithful character adapter identifies every actual
center term with a faithful scalar extension of the corresponding intrinsic
generic block, and the marker is unchanged by that extension.
\end{proof}

\subsection{The projective endpoint and the contradiction}

Formula~\eqref{eq:marker-pbundle} applied to
\(X\times\PP^1=\PP_X(\OO_X^{\oplus2})\) gives
\begin{equation}
 I_{\mathrm{at}}(X\times\PP^1)=2.
 \label{eq:atomic-product-value}
\end{equation}
On the other hand,
\(QH^\bullet(\PP^4)=\Lambda[H]/(H^5-Q)\).  At \(Q\ne0\), multiplication
by \(5H\) has five distinct eigenvalues after algebraic closure.  Thus
\begin{equation}
 I_{\mathrm{at}}(\PP^4)=0.
 \label{eq:atomic-projective-value}
\end{equation}

\begin{proof}[Proof of Theorem~\ref{thm:every-cubic}]
Proposition~\ref{prop:atomic-lowdim} kills every actual center correction in
a weak factorization of smooth projective fourfolds.  Hence
Theorem~\ref{thm:marker-ledger} makes \(I_{\mathrm{at}}\) birationally
invariant.  If \(X\times\PP^1\) were rational, then
\[
 2=I_{\mathrm{at}}(X\times\PP^1)
 =I_{\mathrm{at}}(\PP^4)=0,
\]
a contradiction.
\proves{thm:every-cubic}%
\end{proof}

\begin{remark}[Scope]
The marker, center theorem, and weak-factorization application above are
strictly for one stabilization.  No statement concerns a second or an
arbitrary stabilization.
\end{remark}
